\documentclass[11pt]{scrartcl}

\usepackage[top=1.5cm]{geometry}
\usepackage{url}
\usepackage{hyperref}
\usepackage{float}
\usepackage{listings}
\usepackage{xcolor}

\setlength{\parindent}{0em}
\setlength{\parskip}{0.5em}

\newcommand{\youranswerhere}{[Your answer goes here \ldots]}
\renewcommand{\thesubsection}{\arabic{subsection}}

\lstdefinestyle{dbtsql}{
  language=SQL,
  basicstyle=\small\ttfamily,
  keywordstyle=\color{magenta!75!black},
  stringstyle=\color{green!50!black},
  showspaces=false,
  showstringspaces=false,
  commentstyle=\color{gray}}

\title{
  \textbf{\large Assignment 5} \\
  Concurrency Tuning \\
  {\large Database Tuning}}

\author{
  Group Name (e.g. A1, B5, B3) \\
  \large Lastname1 Firstname1, StudentID1 \\
  \large Lastname2 Firstname2, StudentID2 \\
  \large Lastname3 Firstname3, StudentID3
}

\begin{document}

\maketitle

\paragraph{Notes}

\begin{itemize}
  \item You will need to run transactions concurrently using threads in Java. See \\ \url{https://dbresearch.uni-salzburg.at/teaching/2020ss/dbt/account.zip} for an example.
\end{itemize}

\subsection*{Experimental Setup}

Describe your experimental setup in a few lines.

\youranswerhere{}

\subsection*{Solution (a)}

\paragraph{Read Committed}

Throughput and correctness for solution (a) with isolation level \lstinline[style=dbtsql]{READ COMMITTED}.

\begin{table}[H]
  \centering
  \begin{tabular}{c|c|c}
    \#Concurrent Transactions & Throughput [transactions/sec] & Correctness
      \tabularnewline
    \hline
    1 & \ldots & \ldots \tabularnewline
    2 & \ldots & \ldots \tabularnewline
    3 & \ldots & \ldots \tabularnewline
    4 & \ldots & \ldots \tabularnewline
    5 & \ldots & \ldots \tabularnewline
  \end{tabular}
\end{table}

\paragraph{Serializable}

Throughput and correctness for solution (a) with isolation level \lstinline[style=dbtsql]{SERIALIZABLE}.

\begin{table}[H]
  \centering
  \begin{tabular}{c|c|c}
    \#Concurrent Transactions & Throughput [transactions/sec] & Correctness
      \tabularnewline
    \hline
    1 & \ldots & \ldots \tabularnewline
    2 & \ldots & \ldots \tabularnewline
    3 & \ldots & \ldots \tabularnewline
    4 & \ldots & \ldots \tabularnewline
    5 & \ldots & \ldots \tabularnewline
  \end{tabular}
\end{table}

\subsection*{Solution (b)}

\paragraph{Read Committed}

Throughput and correctness for solution (b) with isolation level \lstinline[style=dbtsql]{READ COMMITTED}.

\begin{table}[H]
  \centering
  \begin{tabular}{c|c|c}
    \#Concurrent Transactions & Throughput [transactions/sec] & Correctness
      \tabularnewline
    \hline
    1 & \ldots & \ldots \tabularnewline
    2 & \ldots & \ldots \tabularnewline
    3 & \ldots & \ldots \tabularnewline
    4 & \ldots & \ldots \tabularnewline
    5 & \ldots & \ldots \tabularnewline
  \end{tabular}
\end{table}

\paragraph{Serializable}

Throughput and correctness for solution (b) with isolation level \lstinline[style=dbtsql]{SERIALIZABLE}.

\begin{table}[H]
  \centering
  \begin{tabular}{c|c|c}
    \#Concurrent Transactions & Throughput [transactions/sec] & Correctness
      \tabularnewline
    \hline
    1 & \ldots & \ldots \tabularnewline
    2 & \ldots & \ldots \tabularnewline
    3 & \ldots & \ldots \tabularnewline
    4 & \ldots & \ldots \tabularnewline
    5 & \ldots & \ldots \tabularnewline
  \end{tabular}
\end{table}

\subsection*{Discussion}

Discuss the outcomes and explain the difference between the isolation levels in PostgreSQL with respect to your experiment.

Explain \textbf{with your own words} how PostgreSQL deals with updates in the different isolation levels, within a transaction and within a single SQL command. Explicitly explain why you got the experimental results for solution (a) and (b), respectively.

\paragraph{Solution (a)}\mbox{}

\youranswerhere{}

\paragraph{Solution (b)}\mbox{}

\youranswerhere{}

\subsection*{Time Spent on this Assignment}

Time in hours per person: \textbf{XXX}

\subsection*{References}

\begin{table}[H]
  \centering
  \begin{tabular}{c}
    \hline
    \textbf{Important:} Reference your information sources! \tabularnewline
    Remove this section if you use footnotes to reference your information sources. \tabularnewline
    \hline
  \end{tabular}
\end{table}

\end{document}
